% Options for packages loaded elsewhere
\PassOptionsToPackage{unicode}{hyperref}
\PassOptionsToPackage{hyphens}{url}
\documentclass[
  11pt,
]{article}
\usepackage{xcolor}
\usepackage[a4paper,margin=2.2cm]{geometry}
\usepackage{amsmath,amssymb}
\setcounter{secnumdepth}{-\maxdimen} % remove section numbering
\usepackage{iftex}
\ifPDFTeX
  \usepackage[T1]{fontenc}
  \usepackage[utf8]{inputenc}
  \usepackage{textcomp} % provide euro and other symbols
\else % if luatex or xetex
  \usepackage{unicode-math} % this also loads fontspec
  \defaultfontfeatures{Scale=MatchLowercase}
  \defaultfontfeatures[\rmfamily]{Ligatures=TeX,Scale=1}
\fi
\usepackage{lmodern}
\ifPDFTeX\else
  % xetex/luatex font selection
\fi
% Use upquote if available, for straight quotes in verbatim environments
\IfFileExists{upquote.sty}{\usepackage{upquote}}{}
\IfFileExists{microtype.sty}{% use microtype if available
  \usepackage[]{microtype}
  \UseMicrotypeSet[protrusion]{basicmath} % disable protrusion for tt fonts
}{}
\makeatletter
\@ifundefined{KOMAClassName}{% if non-KOMA class
  \IfFileExists{parskip.sty}{%
    \usepackage{parskip}
  }{% else
    \setlength{\parindent}{0pt}
    \setlength{\parskip}{6pt plus 2pt minus 1pt}}
}{% if KOMA class
  \KOMAoptions{parskip=half}}
\makeatother
\setlength{\emergencystretch}{3em} % prevent overfull lines
\providecommand{\tightlist}{%
  \setlength{\itemsep}{0pt}\setlength{\parskip}{0pt}}
\usepackage{bookmark}
\IfFileExists{xurl.sty}{\usepackage{xurl}}{} % add URL line breaks if available
\urlstyle{same}
\hypersetup{
  hidelinks,
  pdfcreator={LaTeX via pandoc}}

\author{}
\date{}

\begin{document}

\section{Thesis Defense Transcript}\label{thesis-defense-transcript}

Data-Driven Cognitive Phenotyping in Acquired Brain Injury\\
20-minute presentation script\\
Zoltan Kunos\\
MSc Fundamental Principles of Data Science\\
Universitat de Barcelona\\
Supervisor: Dr.~Alejandro Garcia-Rudolph\\
\textbf{Target duration:} 20 minutes

\section{Slide 1. Title and
Introduction}\label{slide-1.-title-and-introduction}

\textbf{Time on this slide: 45 seconds.}

Hi. Welcome. The title of my dissertation is ``Data Driven Cognitive
Phenotyping in Acquired Brain Injury.''

The principal issue addressed in my dissertation is whether routinely
recorded neuropsychological data can serve as a consistent and reliable
cognitive classification system after acquired brain injury (ABI), even
though the neuropsychological datasets typically are less well-defined
than clinical datasets.

In short, yes. But the cognitive substructures identified by clustering
did not represent clearly defined boundaries separating different
cognitive subgroups. Rather, the most prominent signal detected was a
cognitive severity gradient. Along a general dimension of cognitive
performance, patients were located along this dimension. The clustering
methodology subsequently separated this general dimension of cognitive
performance into three distinct levels of clinically significant
severity.

\textbf{Note:} Provide enough background about the overall argument
before you begin discussing Methods.

\section{Slide 2. Thesis in One
Sentence}\label{slide-2.-thesis-in-one-sentence}

\textbf{Time on this slide: 45 seconds.}

My thesis in one sentence is: A clear cognitive severity gradient can be
identified using unsupervised learning instead of a hidden subgroup
structure in incomplete routine neuropsychological data sets.

There are three key ideas contained in this single sentence. First,
diagnosis does not explain the variability in cognition fully. Two
patients receiving the same high-level diagnosis for ABI may have
differing needs for rehabilitation. Second, missingness must be taken
into account through imputation, as many missing values in the data are
not random missing fields in a table. Typically, they are indicative of
fatigue, cognitive impairment, test intolerance, examiner judgment, or
simply the physical constraints of assessing the degree of cognitive
impairment. Third, no matter which strategy for imputation is chosen,
the tiers that are identified are highly correlated. Therefore, it
increases confidence that the results obtained from my clustering
solution(s) are more credible than those resulting from an arbitrary
preprocessing approach selected without regard to its influence on
clustering. Therefore, my thesis addresses both a clinical and a
methodological question. Clinically, I asked what type of cognitive
organization occurs post-ABI. Methodologically, I asked if that
organization persists notwithstanding the necessity to make decisions
regarding missing data to enable analysis of routinely collected
clinical data.

\section{Slide 3. Clinical Problem}\label{slide-3.-clinical-problem}

\textbf{Time on this slide: 1 min.}

The clinical problem is characterized by heterogeneity. ABI can affect a
wide array of cognitive functions including attention, memory, language,
executive function, visuospatial perception and/orientation to name just
a few. Furthermore, the cognitive impairments exhibited by patients do
not occur in identical manners.

Clinicians rely heavily on broad categories such as diagnosis or
severity of injury to describe patients' cognitive challenges. These
categories provide a valuable conceptual framework for appreciating the
range of potential rehabilitation needs for patients. Nonetheless, these
categories do not possess sufficient detail to allow clinicians to
differentiate rehabilitation needs based on an individual patient's
specific pattern and degree of cognitive difficulty. For example, two
patients exhibiting memory deficits may exhibit substantially disparate
patterns and degrees of cognitive difficulty.

Secondarily, the clinical problem includes missingness. Frequently,
neuropsychological test batteries include substantial gaps as a result
of the reality that neuropsychologists often assess patients during
times when they are least able to successfully complete complex battery
tests. When all incomplete cases are excluded from analysis, essentially
half (approximately 49.9\%) of all assessments made by patients whose
treatment plans would likely benefit most from a comprehensive
description of their cognitive abilities would be removed from
consideration in the analysis.

Therefore, my research question is: Can data-driven clustering generate
clinically meaningful cognitive organizational schemes? Moreover, if
data-driven clustering generates clinically meaningful organizational
schemes, can I trust that such organizational schemes will persist even
when missing data are properly handled rather than being ignored?

\section{Slide 4. Pre-registered
Tests}\label{slide-4.-pre-registered-tests}

\textbf{Time on this slide: 1 minute.}

I developed my analysis plan centered around four quantitative
hypotheses. Hypothesis one (H1) tested whether discrete and stable
clusters existed. To establish whether that was true, I outlined several
pre-specified criteria to evaluate that question. Those criteria
consisted of an average silhouette value \textgreater{} 0.40 and
\textless{} 30\% noise for virtually all of the imputation strategies
examined. Additionally, evidence supporting the existence of at least
two clusters was necessary for the majority of the imputation methods
examined.

Hypothesis two (H2) evaluated the stability of clustering to imputation.
Similar to H1, I specified a threshold for establishing stability. That
is to say that I desired an average pairwise adjusted Rand Index
\textgreater{} .50 across imputation methods. Note once again that H2
utilized a fixed reference design. As such, it employed a single fixed
reference for evaluating similarity - i.e., the embedding and partition
remained unchanged while the values to be imputed were permitted to
vary.

Hypothesis three (H3) tested whether the tiers identified through
clustering possessed clinical significance. If the tiers are clinically
significant, they ought to be related to clinical unit classifications.
I conducted this inquiry using chi-square association statistics,
Cramer's V coefficient and adjustment for diagnosis group.

Lastly, hypothesis four (H4) tested whether aggregating variables at the
domain level resulted in better separation than utilizing all individual
variables. I compared separation metrics, variance inflation factor
coefficients and condition numbers to accomplish this goal.

\textbf{Cue:} Be confident when providing this explanation to the
examiners. Avoid overly detailing statistical formulas here - provide
them simply so that the examiners understand that there is an objective
test of your assertions provided throughout this presentation.

\section{Slide 5. Cohort}\label{slide-5.-cohort}

\textbf{Time on this slide: 1 minute.}

Although my cohort is relatively large, it is incomplete; that is why
this dissertation required both a clustering methodology and a
missing-data strategy.

All analyses were performed using data from a rehabilitation center
neuropsychology database. Following tier-1 filtering and eliminating
assessments that lacked any type of cognitive testing from
consideration, there were 17,406 assessments from 7,285 patients
available for examination. Of those variables assessed in relation to
six cognitive domains, only 14 variables were retained for subsequent
analyses.

Only 53.1 percent of rows in the dataset had complete values across all
14 variables. Complete-case analysis would discard nearly half
(approximately 46.9\%) of all assessments in a clinical rehabilitation
database. Discarding almost half of the assessments in a clinical
rehabilitation database would significantly impact the characteristics
of the patient population analyzed.

Therefore, I view missingness as non-randomly distributed across the
variables assessed; hence, missingness cannot be viewed as an ancillary
characteristic of the datasets analyzed. Missingness reflects various
aspects of the clinical testing process. Assessments completed earlier
in testing may reflect fatigue or decreased motivation; assessments
completed later may reflect cognitive overload or task aversion;
assessments omitted may reflect either insufficient test tolerance or
poor decision-making on behalf of clinicians relative to choosing
alternative forms of assessment; etc.

However, if all incomplete rows are discarded to create a cleaner
numeric representation of the dataset; then that action reduces the
probability that patients with severe cognitive impairment will be
included in analyses examining associations between cognition and
outcomes. Therefore, analyses must utilize imputation strategies
designed to maintain partially observed assessments while allowing for
assessment of uncertainty/sensitivity regarding missing values.

\section{Slide 6. Missingness}\label{slide-6.-missingness}

\textbf{Time on this slide: 1 minute.}

This slide offers additional rationale for addressing missingness as
structured as opposed to random.

Screening tests administered routinely are much more complete than
specialized assessments requiring extended amounts of time to
administer.

As mentioned above, missingness reflects aspects of the clinical testing
process and not simply random error in administration. Examples of
contributing factors to missingness include fatigue,
disorientation/confusion caused by motor deficits or cognitive load or
decreased motivation/task aversion or clinician judgments that another
test may be more appropriate for evaluating certain aspects of a
patient's cognition.

Thus, missingness can affect the ability to define patients'
rehabilitation requirements in terms extending beyond simply diminishing
sample size. Removing partially observed rows will also tend to remove
patients with higher degrees of cognitive deficit.

Therefore, analyses must employ imputation methodologies enabling
retention of partially observed assessments while permitting evaluations
of uncertainty/sensitivity relating to missing values.

\section{Slide 7. Pipeline}\label{slide-7.-pipeline}

\textbf{Time on this slide: 1 minute and 20 seconds.}

Below is illustrated my overall pipeline flow chart.

Illustrated in the flow chart is where I employed imputation
methodologies to manage missingness in my datasets. Ten different
imputation strategies ranging from simple/basic to
model-based/donor-based/matrix completion/machine learning/deep learning
approaches were evaluated.

After employing imputation methodologies, I normalized/stabilized scores
across the imputed datasets and categorized/grouped them by six
cognitive domains. The imputed datasets were subsequently embedded using
Uniform Manifold Approximation and Projection (UMAP). Finally,
hierarchical density-based spatial clustering of applications with noise
(HDBSCAN) was used to determine clusters in embedded space. Separation
vs percentage of noise in each cluster configuration evaluated was
determined by calculating silhouette scores as a quality measure.

For hypothesis two (H2), I employed a fixed reference design. Utilizing
Multiple Imputation By Chained Equations (MICE), I designated an anchor
scaler/embedding/HDBSCAN cluster configuration (fixed reference) to
project other imputation methods through this fixed structure.
Importantly for H2, I tested whether imputed values reside in similar
locations given a fixed structure versus separately fitting a full UMAP
+ HDBSCAN pipeline produces equivalent hierarchies

\section{Slide 8. H1 results}\label{slide-8.-h1-results}

\textbf{Estimated duration: 1 min.}

three cognitive severity tiers are evident. Labels for tiers were
defined as above average (aa); near normal (nn), and Global Impairment
(gi). Aa contained 6,725 assessments. Patients scoring higher than the
cohort mean across domains. Nn contained 6,328 assessments. Situated
near the cohort average. Gi contained 4,353 assessments. Below the
cohort mean across domains.

All 10 Imputation methods exceeded the pre-registered thresholds. For
the mice reference solution, the full coverage silhouette was about .53.
The recovered number of tiers was three. I am unable to understand these
as natural kinds or sharp distinctions between phenotypes. Instead, they
appear as ordered bands along a severity continuum.

\section{Slide 9. Severity profile}\label{slide-9.-severity-profile}

\textbf{Approximately 1 min.}

first major finding: a significant cognitive severity gradient exists
following acquired brain injury (abi)

Principal components analysis (pca): confirms trend in data. Pca
explains nearly 56 percent of variance in data and all domains have
positive loads. Loading pattern consistent with broad cognitive severity
dimension

Most prominent signal in data not symptom specific dissociation or
catalog of distinct subtypes. Most dominant signal is overall cognitive
severity

\section{Slide 10. H2 results}\label{slide-10.-h2-results}

\textbf{Time: approximately 1 minutes and 15 seconds.}

Second major finding relates to robustness of imputed values

Each pair of methods from each comparison met predetermined criteria for
agreement established prior to conducting analyses.

Mean adjusted Rand index (ARI) among pairs of methods was .71, based
upon series of bootstraps with a 95\% confidence interval (.70,.72).
Thus indicated substantial agreement among method pairs. Of the 45
comparisons made to assess agreement, 44 survived holm's correction.
Only one pair did not meet predetermined threshold. That pair was KNN
vs.~SoftImpute, which had a pairwise ARI just short of predetermined
threshold.

High-confidence consensus labels were assigned to more than 86\% of all
assessments. This suggests little variation between selected Imputation
methods relative to fixed embedding and partition. However, scope of
claim only extends to demonstrate point assignments are relatively
robust under fixed embedding and partition not independent clustering
pipelines fitted to different datasets will produce identical density
hierarchies.

\section{Slide 11. Why Imputation
matters}\label{slide-11.-why-imputation-matters}

\textbf{Time: approximately 1 minute.}

Illustrated necessity for Imputation analysis. Complete cases represent
only 53.1 percent of original 17,406 assessments.

When evaluating similarity between mice reference solution and baseline
complete case solution, both solutions produced two clusters rather than
desired three-tiered solution. Silhouette values also indicate baseline
solution produces poorer silhouette values (about .26) compared to those
produced by mice reference solution (about .53). Lastly, although ARI
comparing the two solutions on common assessments was about .63, still
represents less than perfect agreement.

Thus deleting cases due to missing data does not merely produce
less-powerful version of same analysis. It changes partition. How
missing data are handled forms part of model's output regarding its
substantive outcome.

\section{Slide 12. H3 result}\label{slide-12.-h3-result}

\textbf{Time: approximately 1 minute.}

third hypothesis examines whether cognitive severity levels identified
in previous slides relate to clinical unit assignment. Chi-square test
between tier membership and clinical unit assignment yields a chi-square
statistic equal to 911 and estimated p value greater than 10\^{}-156.
Cramer's v = .16 indicates small-to-moderate effect size according to
standard interpretations. As per pre-registration document, designated
an effect size of .10 as a threshold. Since cramer's v = .16
\textgreater{} .10 consider this result supporting h3.

Also employed cochran-mantel-haenszel test to determine if clinical unit
assignment can be solely explained by diagnoses. Even after controlling
for diagnosis in cmh test, relationship remained statistically
significant. Conclude while diagnostic classification plays role in
determining clinical unit assignment, cognitive severity tier provides
additional predictive power above-and-beyond diagnostic classification.

From a clinical perspective reasonable to assume individuals classified
as having Global Impairment will require more intensive structured
rehabilitative interventions than their peers who are not classified in
this manner. Do not mean tier membership predetermines placement
decisions in rehabilitation settings. Propose tier membership
contributes valuable information to inform rehabilitation routing
decisions.

\section{Slide 13. H4 result}\label{slide-13.-h4-result}

\textbf{Time: approximately 1 minute.}

Final hypothesis assesses if cluster solutions derived from aggregating
domain-level scores perform comparable to cluster solutions generated
using unaggregated (ie individual test variable) scores.

Domain level clustering performs marginally better than unaggregated
individual test variable clustering. Silhouette values indicate marginal
improvement in clustering quality (.48 for domains vs.~.47 for
Individual Test Variables). Additionally, when evaluating numerical
stability via variance inflation factors (vifs), mean vif decreases
significantly (from \textasciitilde2.80 to \textasciitilde1.90) and
condition numbers improve dramatically (from \textasciitilde62 to
\textasciitilde12).

Improvements are especially relevant for downstream applications
involving distance-based clustering algorithms (eg UMAP, hdbscan), given
such algorithms often highly sensitized to redundancy/correlation
between features. By aggregating individual test scores into clinically
meaningful domains, reduced feature redundancy and retained
interpretable cognitive constructs.

As such, domain level clustering offers advantages in terms of
computational stability and interpretability making it preferable to
unaggregated clustering approaches when attempting to identify
cognitively-distinctive patient populations.

\section{Slide 14. Result Summary}\label{slide-14.-result-summary}

\textbf{Duration of this slide: approximately 45 seconds.}

Summary of four tested hypotheses

Hypothesis \#1 finds support: recovered stable clusters using
unsuptervised techniques but understand these as representing distinct
levels of cognitive severity rather than categorically-distinct
``subtypes.''

Hypothesis \#2 finds support within context of fixed-reference design:
average pairwise ARI \textasciitilde{} .71 suggests robust agreement
among method pairs.

Hypothesis \#3 finds support: severity tier related to clinical unit
assignment and continued statistical significance even after controlling
for diagnosis.

Hypothesis \#4 finds support: domain level clustering exhibits superior
numerical properties (silhouette \& vif) compared to unaggregated test
variable clustering.

Collectively findings provide empirical support for developing and
applying missingness aware severity stratification approach for
analyzing routine neuropsychological assessments collected post abi.

\section{Slide 15. Interpretation}\label{slide-15.-interpretation}

\textbf{Duration: approximately 1 min., 15 sec.}

Primary interpretation offered here is continuum of cognitive severity
gradients existent in dataset; rather than collection of discrete
subtypes.

Understanding this distinction is important. An interpretation of
subtypes suggests qualitative form(s) of cognitive profile(s) exist. On
the otherhand, interpretation of phenomenon as example of cognitive
severity gradient suggests differences in cognitive function lie almost
exclusively along an overall continuum of cognitive severity across
multiple domains.

Evidence provided throughout manuscript supports latter viewpoint.
Evidence included radial plots which are indicative of concentricity
among all groups examined via radar plots; large proportion
(\textasciitilde50\%) of total variance explained by pc1; all domains
exhibit positive loads on pc1. Further evidence supporting view comes
from examining conceptual nature of each tier examined during current
investigation which can be interpreted analogously to bands extracted
from larger continuum.

While it may be possible that domain-specific dissociations exist
following brain injury (eg selective deficits in memory following tbi),
evidence presented in current manuscript suggests that pipeline
utilizing Imputation-robust domain-level severity stratified framework
yields stronger signals regarding overall cognitive severity across
multiple domains than any domain-specific dissociation.

Finally, evidence presented here highlights potential pitfalls in
identifying patterns resembling phenotypic differentiation spurred by
poor scaling (and/or inappropriate scaling), use of non-cognitive
measures, omission of missing data.

\section{Slide 16. Utility - Clinical and
Analytical}\label{slide-16.-utility---clinical-and-analytical}

\textbf{Approximately 1 Minute}

Two potential groups of people could utilize this methodology and
subsequent applications:

Clinical Users. The severity tiers could also be utilized as another
source of support for clinicians when making the determination of
rehabilitation intensity, and possibly help clinicians report cognitive
assessment results along with associated confidence levels that can help
alleviate misinterpretation of borderline cases.

Analytical Users. The present manuscript illustrates an example
missing-data-aware pipeline that can be reproduced and modified by
researchers looking to examine relationships between clinical outcomes
and cognitive severity following ABI.

Ultimately, the end goal of this dissertation will be to increase
transparency about how missing data are treated within statistical
analyses. As opposed to trying to ignore missing data or eliminate all
incomplete observations from analyses, researchers will now be able to
make missing-data decisions transparent and quantifiable in addition to
evaluating the impact they have on the resulting analyses.

\section{Slide 17. Application:
CogDash}\label{slide-17.-application-cogdash}

\textbf{Approximately 45 seconds}

In order to improve usability, I developed CogDash - a Streamlit web-app
that utilizes portions of my analyses to create an interactive
decision-making tool for clinicians.

The CogDash app contains tools such as cluster explorers for each
patient; patient look-up; classification of new patients into clusters;
summaries of tested hypotheses; and access to advanced analytical tools.
Overall, CogDash was designed to promote collaboration/understanding
between clinicians wanting to assess the cognitive function of their own
patients, and analysts providing analytical services; it is NOT meant to
automatically determine clinical decisions.

One possible reason why CogDash has value is because models applied in
rehabilitation environments need to be auditable. Clinicians need to be
able to view the data used and at what points confidence is low.

\section{Slide 18. Limitations}\label{slide-18.-limitations}

\textbf{Approximately 1 Min., 15 Sec.}

There are many limitations to this study:

Firstly, data was collected from one location, so replication studies
are required before we consider this structural representation useful
across locations collecting post-ABI neuropsychological data.

Secondly, this analysis did not represent a complete longitudinal
recovery model. Although repeated assessments provided some information
on recovery trajectories, including time since injury added additional
context; and direct modeling of recovery would require additional
longitudinal designs.

Thirdly, the Imputation strategy rested heavily on assumptions related
to MAR missing-data-generating processes; however, it cannot be ruled
out that MNAR processes exist based solely on observable data.
Specifically, extreme MNAR analyses indicated that global impairment
tier is robust while the boundaries around that tier are likely subject
to MNAR processes.

Finally, there were no available functional outcomes (e.g.,
return-to-work; LOS) to assess for validation purposes; therefore, these
types of functional outcomes would provide the most valid external
validators for future research examining whether tier membership
provides utility for guiding rehabilitation plans.

\section{Slide 19. Conclusion}\label{slide-19.-conclusion}

\textbf{Approximately 1 Minute}

I will summarize three findings:

Firstly, routinely-collected post-ABI clinical neuropsychological data
collect significant amounts of cognitive-severity gradient.

Secondly, the method of imputing missing data represents substantial
modeling decisions and affect the resultant structure(s), therefore
ignoring missing data as part of a minor preprocessing procedure could
lead to different results depending on characteristics of available
data.

Thirdly, severity-tiered representations by domain level represent
clinically-usable \& numerically-stable representations representing
benefits compared to un-aggregated clustering approaches when
identifying cognitively-distinctive patient populations.

Therefore collectively, this dissertation supports neither strong
evidence for distinct cognitive phenotypes existing in this sample
population; nor does it negate their existence completely; however it
does support developing \& applying an Imputation-robust
severity-stratification framework that can be both
clinically-interpretable and analytically-evaluated.

\section{Slide 20. Thank You!}\label{slide-20.-thank-you}

\textbf{Approximately 30 Seconds}

Thank you for your time. I'd like to answer any questions you may have
pertaining to either utilizing my codebase/Dashboards, clarifying
technical details (Imputation design/clustering choices/etc.), or
discussing potential implications for future research projects examining
cognitive-function/rehabilitation planning post-ABI.

\textbf{Optional Final Sentence if Room is Silent}: Overall, my hope is
that the analyses shown here identified something
clinically-useful-but-more conservatively-defined as a
cognitive-severity continuum after acquired-brain-injuries rather than
distinct cognitive subtypes

\end{document}
